% Part: incompleteness
% Chapter: theories-computability
% Section: introduction

\documentclass[../../../include/open-logic-section]{subfiles}

\begin{document}

\olfileid[fr]{inc}{tcp}{int}

\olsection{Introduction}

Ce chapitre étudie la calculabilité d'ensembles d'énoncés
à partir de leur codage arithmétique.

Nous disposons maintenant des résultats suivants :
\begin{enumerate}
\item une définition de la représentabilité d'une fonction totale
  dans~$\Th{Q}$
  (\olref[req][int]{defn:representable-fn});
\item une définition de la représentabilité d'une relation
  dans~$\Th{Q}$
  (\olref[req][rel]{defn:representing-relations});
\item un théorème selon lequel les fonctions totales représentables
  dans~$\Th{Q}$ sont exactement les fonctions totales calculables
  (\olref[req][int]{thm:representable-iff-comp});
\item un théorème selon lequel les relations représentables
  dans~$\Th{Q}$ sont exactement les relations calculables
  (\olref[req][rel]{thm:representing-rels}).
\end{enumerate}

Une {\em théorie} est un ensemble d'!!{sentence}s clos par
déduction : si $T$ démontre~$!A$, alors $!A$ appartient à~$T$.
On peut voir une théorie comme un ensemble d'!!{sentence}s
auquel on a ajouté toutes leurs conséquences.
Désormais, $\Th{Q}$ désignera la {\em théorie} constituée
de tous les !!{sentence}s déductibles des huit axiomes
donnés en \olref[req][int]{sec}.
Rappelons que nous pouvons coder les !!{formula}s de
$\Th{Q}$ par des nombres; si $!A$ est une telle
!!{formula}, $\Gn{!A}$ désigne son numéro de Gödel.
Grâce à ce codage, nous pouvons demander si divers
ensembles de !!{formula}s sont calculables.

\end{document}
